\section{Discussion}

The main value of City1 v3 is not that it turns a proxy surface into truth. Its value is that it reduces deterministic over-reading. A calibrated proxy map can otherwise look uniformly authoritative even when the supporting evidence is clearly uneven. By adding intervals, confidence bands, and hotspot classes, v3 gives the user a more disciplined way to separate stronger and weaker screening cases.

Almaty is the clearest example of why this matters. It combines the strongest freeze-city interpretation note, the largest high-confidence share, and the largest number of stable high-value hotspot cells. In practical terms, this means Almaty can support more assertive exploratory screening than the other frozen cities, while still remaining inside the proxy framework. The paper does not claim that these cells are true census hotspots. It claims that they are the strongest screening candidates in the frozen evidence package.

Astana and Shymkent illustrate the opposite side of the story. Their city-level outputs are calibrated and still useful, but the confidence distribution and hotspot composition are caution-heavier. In both cities, the uncertainty layer acts less like an endorsement and more like a warning against over-interpretation. This is an important scientific contribution in its own right: a useful uncertainty system should sometimes tell the analyst to slow down rather than only decorate the output.

Semey is the intermediate case. It is neither as strong as Almaty nor as caution-heavy as Shymkent. It retains a small but real stable hotspot subset together with a larger review-oriented tail. That makes it a useful example of how v3 can support screening in a mixed-support city without pretending that support is uniformly strong.

The strongest support for v3 comes from hotspot stability and uncertainty-burden separation, not from every validation layer equally. Stable classes consistently retain much higher mean stability metrics than caution-heavy ones, and confidence bands separate relative uncertainty in the expected direction. By contrast, interval coverage is mixed, external disagreement alignment is mixed or negative, and district interval coverage remains partial. These weaker layers constrain the claim, but they do not nullify the stronger evidence. They simply narrow the correct interpretation: v3 is an uncertainty-aware screening layer, not a full uncertainty-calibrated truth model.

This bounded interpretation also makes the next project step clearer. Future tool-grounded interpretation layers, including planned v4 work, should consume confidence-aware outputs rather than only a deterministic median surface. But that future layer should sit on top of the current v3 evidence stack, not be confused with the method presented here.

